\chapter{Project Summary}
\sffamily
\label{chap:Project Summary}
%%%%%%%%%%%%%%%%%%%%%%%%%%%%%%%%%%%%
The Aerodynamic Test Platform and Wings (ATPW) team in the 2025-2026 year consolidates three large aspects of the overall Micro Flapping Wing Flyer (MFWF) project. The first focuses on the manufacturing and operation of a test platform for which the MFWF is to use to measure the orthogonal forces it produces under self-actuation. The previous capstone year succeeded in developing a prototype of this test platform but were plagued by hardware and shipping constraints, resulting in an insufficient amount of testing and finalization of the platform. This year, the platform went through significant testing and small redesigns to validate its functionality and reliability in terms of measuring lift-averaged forces of a flapping actuator. The result of this year is a commissioned instrument that can qualitatively discern differences in lift production of an actuator.

The second goal of the ATPW team involved developing and training a neural network model that could predict forces based on experimental inputs. This progressed over the Fall 2025 and Winter 2026 semesters. The Fall semester included the development and validation of an initial proof-of-concept aerodynamic prediction model using machine learning and structured, static, NACA airfoil data. The work in the Winter semester shifted towards a more physically representative prediction model of unsteady three-dimensional aerodynamic forces and moments. The overall results of this year established a foundation for future training using Aerodynamic Test Platform (ATP) and Rotational Rates Test Stand (RRTS) experimental data as it becomes available.

The final purpose of the ATPW team this year was to construct a strong foundation towards defining and characterizing the MFWF wing structure capable of flight. The previous capstone year had performed static testing on a large variety of vein structures, but failed to capture any dynamic results. This year, the goal was to define a DOE to understand how every aspect of wing structure design can affect the motion and ultimately aerodynamic effectiveness of the wing. A full test procedure was developed to quantitatively define the optimal wing upon further testing. This will in future, result in a high performance flapper able to produce sufficient lift and accurate movement.

Throughout the 2025-2026 academic year, the ATPW team has made significant advances in all three of these aspects. The following is a summary of the major advances that have been made in each of these three areas.
\begin{enumerate}
    \item Aerodynamic Test Platform
    \begin{enumerate}
        \item Iterative redesigns of the platform's structure and electronics to improve the resolution of the force measured
        \item Standardization of actuator mounting methods to the test platform
        \item Validation and analysis of the platform's ability to measure lift-averaged forces
    \end{enumerate}
    \item Neural Network Development
    \begin{enumerate}
        \item Gain understanding of neural networks
        \item Develop a proof-of-concept machine learning application using static aerodynamic data
        \item Develop a machine learning application using a dynamic dataset comparable to anticipated experimental results from the ATP and RRTS.
    \end{enumerate}
    \item Wing Design and Testing 
    \begin{enumerate}
        \item Streamlined design process for ease of manufacturing.
        \item Review of published literature on flapping Wing designs and testing.
        \item Iterative design and manufacturing of wings
        \item Design and commission a dynamic test stand with high speed camera imaging.
    \end{enumerate}
\end{enumerate}

